\section{GIỚI THIỆU ĐỀ TÀI}

\subsection{Bối cảnh và động lực}
Trong kỷ nguyên số hóa hiện nay, cuộc chạy đua vũ trang giữa các chuyên gia bảo mật và tội phạm mạng đang diễn ra với tốc độ chưa từng có. Các hệ thống phát hiện xâm nhập (Intrusion Detection Systems - IDS) và phần mềm diệt virus truyền thống, vốn dựa trên chữ ký (signature-based), đang ngày càng trở nên lỗi thời trước sự gia tăng của mã độc đa hình (polymorphic malware) và các cuộc tấn công zero-day. Sự ra đời của Trí tuệ Nhân tạo (AI) và Học máy (Machine Learning - ML) đã mang lại một bước tiến lớn, cho phép phát hiện mã độc dựa trên hành vi và các đặc trưng tĩnh/động phức tạp thay vì chỉ so sánh chuỗi byte đơn thuần. Tuy nhiên, mô hình huấn luyện tập trung (Centralized Learning) - nơi dữ liệu từ hàng triệu thiết bị đầu cuối được gửi về một trung tâm dữ liệu khổng lồ để xử lý - đang vấp phải những rào cản kỹ thuật và pháp lý nghiêm trọng. \\

Thứ nhất, vấn đề quyền riêng tư dữ liệu đang trở thành tâm điểm với sự ra đời của các quy định như GDPR (Châu Âu) hay CCPA (California). Dữ liệu nhật ký mạng, file thực thi, hoặc hành vi người dùng trên thiết bị di động chứa đựng thông tin nhạy cảm mà các tổ chức không muốn hoặc không được phép chia sẻ ra bên ngoài. Thứ hai, chi phí băng thông và độ trễ mạng khi truyền tải hàng terabyte dữ liệu thô về máy chủ trung tâm là một gánh nặng khổng lồ đối với cơ sở hạ tầng mạng.\\

Trong bối cảnh đó, Học liên kết (Federated Learning - FL) nổi lên như một giải pháp kiến trúc mang tính cách mạng. Bằng cách đảo ngược quy trình học tập - di chuyển tính toán đến nơi có dữ liệu thay vì di chuyển dữ liệu đến nơi tính toán - FL giải quyết đồng thời bài toán bảo mật và hiệu năng. Báo cáo này sẽ trình bày một nghiên cứu toàn diện và chi tiết về việc thiết kế, triển khai và giám sát một hệ thống phát hiện mã độc phân tán sử dụng Framework Flower (flwr) \cite{beutel2022flowerfriendlyfederatedlearning}, kết hợp với giao diện giám sát thời gian thực dựa trên web. Chúng ta sẽ đi sâu vào từng lớp của kiến trúc, từ xử lý dữ liệu đặc trưng bộ nhớ MalMem2022 (memory dump) tại biên, xây dựng mô hình Neural Networks, tùy chỉnh chiến lược tổng hợp (Aggregation Strategy) trong Flower \cite{beutel2022flowerfriendlyfederatedlearning}, đến việc trực quan hóa luồng huấn luyện thông qua giao diện giám sát thời gian thực. \\

Song song với việc triển khai FL, việc ứng dụng Quantum Machine Learning (QML) trong các bài toán phân loại (classification) hiện cũng đang là một hướng nghiên cứu tiên tiến được quan tâm rộng rãi . Do đó, chúng tôi quyết định khảo sát lý thuyết ba kiến trúc gồm Quantum Convolutional Neural Network \cite{Cong_2019}, Quantum Principal Component Analysis \cite{Lloyd_2014}, và Quantum Multilayer Perceptron \cite{shao2018quantummodelmultilayerperceptron}, đồng thời hiện thực hóa thực nghiệm một mô hình Quantum Neural Network lai (Hybrid QNN) với 4 qubit để kiểm chứng tính khả thi.

Việc tích hợp Quantum Machine Learning được kỳ vọng sẽ khắc phục một số hạn chế còn tồn tại của Federated Learning truyền thống, cụ thể như: khả năng bảo mật dữ liệu vượt trội nhờ các nguyên lý cơ học lượng tử, giảm thiểu chi phí truyền thông do số lượng tham số mô hình nhỏ hơn, và khả năng biểu diễn tốt hơn đối với các cấu trúc dữ liệu phức tạp. Tuy nhiên, do những hạn chế về tài nguyên phần cứng hiện tại, nghiên cứu này sẽ giới hạn phạm vi ở việc đánh giá hiệu năng trên số lượng qubit nhỏ nhằm mục đích kiểm chứng tính khả thi của phương pháp đề xuất.

\subsection{Phương pháp nghiên cứu}

Để đảm bảo tính khoa học và hiệu quả của hệ thống, phương pháp nghiên cứu được chia thành các bước cụ thể sau:

Đầu tiên, để lựa chọn mô hình tối ưu cho máy chủ, chúng tôi quyết định thực hiện benchmark trên nhiều ứng viên mô hình khác nhau nhằm tìm ra sự cân bằng giữa tốc độ huấn luyện và độ chính xác dự báo.

Tiếp theo, nhóm nghiên cứu tiến hành xây dựng cơ chế aggregation, thiết lập máy chủ Flower và triển khai các giải pháp bảo mật cho kênh truyền thông giữa máy chủ và các thiết bị biên. Sau đó, một dashboard trực quan hóa dữ liệu được xây dựng để theo dõi và trình bày kết quả phân tích theo thời gian thực.

Cuối cùng, nghiên cứu sẽ tập trung vào việc đánh giá tác động thực nghiệm của việc tích hợp Quantum Machine Learning đối với quy trình Federated Learning, từ đó rút ra các kết luận về tiềm năng và hạn chế của phương pháp lai này.